% !TeX encoding = UTF-8
\documentclass[variant=apuntes,cover=institutional,mode=screen]{ua-engineering-v6-article}
% !TeX encoding = UTF-8
\UASetUniversity{Universidad Austral}
\UASetFaculty{Facultad de Ingeniería}
\UASetDepartment{Departamento de Computación}
\UASetCampus{Pilar}
\UASetCareer{Ingeniería Informática}
\UASetCommission{Comisión A}
\UASetProfessor{Equipo docente}
\UASetSemester{Segundo cuatrimestre 2026}
\UASetCourse{Sistemas Distribuidos}
\UASetDate{2026}


\UASetClassNumber{12}
\UASetClassTitle{Chaos Engineering y validación bajo falla}
\UASetDocKind{Apuntes}

\title{Clase 12 --- Chaos Engineering, fault injection y testing distribuido}
\author{\UAProfessor}
\date{\UADate}

\begin{document}
\maketitle

\section{Introducción}

En sistemas distribuidos, diseñar mecanismos de resiliencia no basta. El verdadero problema es demostrar que esos mecanismos se comportan como se espera bajo condiciones de falla realistas.

Una arquitectura puede parecer correcta en papel y aún así colapsar cuando:
\begin{itemize}
\item la latencia sube un poco;
\item una dependencia responde intermitentemente;
\item una cola se llena;
\item un retry mal calibrado amplifica la presión;
\item una partición parcial rompe una hipótesis implícita.
\end{itemize}

Por eso, la validación experimental del comportamiento bajo falla es una disciplina central.

\section{Testing distribuido y sus límites}

\subsection{Testing clásico}

Unit tests e integration tests tradicionales son necesarios, pero tienen límites evidentes:
\begin{itemize}
\item suelen asumir tiempos ideales;
\item usan mocks demasiado limpios;
\item rara vez modelan pérdida, reorder o partición;
\item prueban funciones, no siempre propiedades sistémicas.
\end{itemize}

\subsection{Problema}

Los errores distribuidos aparecen precisamente en:
\begin{itemize}
\item interacciones;
\item concurrencia;
\item sincronización temporal;
\item propagación de falla;
\item reacciones colectivas del sistema.
\end{itemize}

\begin{uawarn}
Un sistema puede pasar todos sus tests clásicos y aun así fallar gravemente ante perturbaciones realistas.
\end{uawarn}

\section{Chaos Engineering}

\subsection{Definición}

Chaos Engineering es la disciplina experimental que estudia el comportamiento del sistema bajo perturbaciones controladas.

No consiste en “romper cosas” indiscriminadamente, sino en ejecutar experimentos deliberados con hipótesis explícitas.

\subsection{Estructura de un experimento}

Un experimento de chaos razonable tiene:

\begin{enumerate}
\item una definición de \emph{steady state};
\item una hipótesis de comportamiento;
\item una perturbación controlada;
\item una observación medible del resultado;
\item una conclusión que permita aprender o corregir.
\end{enumerate}

\section{Steady State}

\subsection{Definición}

El steady state es un conjunto de señales que representan comportamiento aceptable del sistema.

Ejemplos:
\begin{itemize}
\item tasa de requests exitosas;
\item latencia p95;
\item throughput sostenido;
\item longitud de cola;
\item porcentaje de errores;
\item tiempo de propagación de eventos.
\end{itemize}

\subsection{Importancia}

Sin steady state, el experimento no tiene criterio de éxito o fracaso. Solo introduce perturbación, pero no produce conocimiento verificable.

\section{Fault Injection}

\subsection{Idea}

Fault injection consiste en introducir fallas o degradaciones controladas para observar la respuesta del sistema.

\subsection{Tipos frecuentes}

\begin{itemize}
\item crash de proceso;
\item reinicio abrupto;
\item aumento artificial de latencia;
\item pérdida de paquetes;
\item throttling de CPU;
\item saturación de disco;
\item errores de dependencia;
\item particiones de red;
\item colas artificialmente lentas.
\end{itemize}

\subsection{Objetivo}

El objetivo no es verificar que el sistema “nunca falla”, sino observar:
\begin{itemize}
\item si degrada con gracia;
\item si los mecanismos de resiliencia se activan;
\item si el daño queda contenido;
\item si la observabilidad permite entender qué pasó.
\end{itemize}

\section{Qué propiedades se validan}

Una buena campaña de fault injection puede validar:

\begin{itemize}
\item timeouts correctamente calibrados;
\item retries acotados;
\item circuit breakers efectivos;
\item backpressure funcional;
\item bulkheads realmente aislantes;
\item compensaciones correctas;
\item recuperación razonable del steady state.
\end{itemize}

\section{Testing distribuido por capas}

\subsection{Unit tests}
Validan lógica local. Son insuficientes para propiedades distribuidas.

\subsection{Integration tests}
Validan integración de pocos componentes. Útiles, pero normalmente demasiado acotados.

\subsection{System tests}
Validan comportamiento del sistema completo en escenarios nominales.

\subsection{Failure tests}
Inyectan fallas específicas y verifican respuesta.

\subsection{Chaos experiments}
Exploran sistemáticamente perturbaciones bajo hipótesis operacionales.

\section{Hipótesis experimentales}

Una hipótesis útil tiene forma concreta, por ejemplo:

\begin{quote}
Si el servicio de pagos responde con latencia alta, el breaker debe abrir antes de que el checkout completo se vuelva indisponible.
\end{quote}

Esto es mejor que una formulación vaga como:
\begin{quote}
“queremos ver si aguanta.”
\end{quote}

\section{Blast Radius}

\subsection{Definición}

El blast radius es el alcance potencial del daño que puede causar el experimento.

\subsection{Importancia}

Todo experimento debe controlar explícitamente:
\begin{itemize}
\item qué componente se afecta;
\item qué porcentaje de tráfico;
\item qué servicios vecinos pueden verse afectados;
\item cómo abortar rápidamente.
\end{itemize}

\section{Seguridad operacional}

Chaos Engineering responsable exige:

\begin{itemize}
\item observabilidad suficiente;
\item límites claros de abortar;
\item entorno adecuado;
\item experimentos pequeños antes que grandes;
\item objetivos explícitos;
\item comunicación con los equipos involucrados.
\end{itemize}

\subsection{Nunca hacer}

\begin{itemize}
\item experimentos sin hipótesis;
\item experimentos sin medir steady state;
\item experiments de alto impacto sin rollback;
\item perturbaciones masivas en producción sin control fino.
\end{itemize}

\section{Ejemplo razonado}

Consideremos un sistema de checkout.

Hipótesis:
\begin{quote}
Si Payments sufre latencia alta, el sistema de checkout degrada parcialmente pero no colapsa.
\end{quote}

Perturbación:
\begin{itemize}
\item agregar 800 ms de latencia a Payments.
\end{itemize}

Señales esperadas:
\begin{itemize}
\item sube la latencia de checkout;
\item el breaker de pagos se activa;
\item se mantiene estabilidad en otras rutas;
\item no aparece explosión de retries;
\item el sitio no se cae por completo.
\end{itemize}

Este experimento no prueba “todo el sistema”, pero sí valida una propiedad operacional concreta.

\section{Relación con observabilidad}

Chaos Engineering depende profundamente de la observabilidad:
\begin{itemize}
\item sin métricas no hay steady state;
\item sin trazas no se entiende la propagación;
\item sin logs no se reconstruyen eventos específicos.
\end{itemize}

\begin{uakeybox}
Observabilidad permite ver la falla. Chaos Engineering permite aprender deliberadamente de ella.
\end{uakeybox}

\section{Modelo del curso}

Podemos expresar el problema como:

\[
D = (P, F, G, S, T)
\]

Por ejemplo:
\begin{itemize}
\item $P$: validar resiliencia real de un sistema distribuido;
\item $F$: latencia, crash, pérdida, saturación, partición;
\item $G$: mantener steady state aceptable o degradación contenida;
\item $S$: fault injection + experimentación controlada + observabilidad;
\item $T$: complejidad operativa, riesgo experimental y costo de instrumentación.
\end{itemize}

\section{Conclusión}

Chaos Engineering y fault injection permiten transformar la resiliencia desde una afirmación teórica hacia una propiedad validada empíricamente.

\begin{quote}
No alcanza con creer que el sistema resistirá. Hay que demostrar, bajo perturbación realista, qué tan bien falla.
\end{quote}

\begin{uakeybox}
La disciplina no consiste en producir caos, sino en producir conocimiento confiable sobre el comportamiento del sistema bajo falla.
\end{uakeybox}


\section{Síntesis razonada de la clase}
Esta unidad debe leerse como parte de una cadena de decisiones de diseño y no como un catálogo de mecanismos. Para estudiar el tema con profundidad conviene reconstruir cuatro preguntas: \emph{qué problema aparece}, \emph{qué garantía se desea}, \emph{qué mecanismo la aproxima} y \emph{qué costo introduce}. En cada ejemplo debe identificarse además el modelo de falla implícito.

\begin{uakeybox}
Una respuesta técnicamente madura no termina al nombrar un algoritmo o patrón: declara sus supuestos, la propiedad que preserva y el escenario en el que deja de ser suficiente.
\end{uakeybox}

\section{Bibliografía guiada y ruta de profundización}
Para profundizar esta clase se recomienda: Principles of Chaos Engineering; Google SRE; Jepsen como referencia de validación.

La lectura debe hacerse con preguntas concretas: (1) ¿qué modelo de sistema asume el autor?, (2) ¿qué propiedad demuestra o persigue?, (3) ¿qué falla o anomalía motiva el mecanismo?, y (4) ¿qué trade-off queda abierto?

\section{Conexión con el Trabajo Final Integrador}
El grupo debe señalar qué parte de su sistema utiliza los conceptos de esta unidad, justificar por qué son necesarios y documentar una alternativa descartada. Cuando el contenido todavía no corresponda a la etapa vigente, debe registrarse como hipótesis o decisión pendiente.

\section{Autoevaluación conceptual}
\begin{enumerate}
\item ¿Cuál es el problema distribuido concreto que motiva esta clase?
\item ¿Qué propiedad de seguridad y qué propiedad de progreso están en juego?
\item ¿Qué supuesto de red, tiempo o falla resulta decisivo?
\item ¿Qué trade-off aceptaría en un sistema real y cuál no aceptaría en un dominio crítico?
\item ¿Cómo observaría experimentalmente que la solución se comporta como espera?
\end{enumerate}

\section{Profundización autocontenida v9}
\subsection{Problema, límite e invariante}
El núcleo de esta clase es testing distribuido, fault injection, steady state, hipótesis, blast radius y diseño experimental. Para estudiarlo de manera autónoma conviene separar tres planos. Primero, el \emph{problema observable}: qué comportamiento incorrecto puede ver un cliente o un operador. Segundo, el \emph{límite del modelo}: qué información, sincronía o coordinación no está disponible. Tercero, el \emph{invariante}: qué propiedad no debe violarse aunque el sistema pierda progreso temporalmente.

El modelo de fallas relevante incluye perturbación sin hipótesis, blast radius excesivo, señal insuficiente y falso positivo. Una solución debe describirse como protocolo y no solo como nombre de tecnología: actores, estados, mensajes, condición de éxito, condición de aborto y estado visible después de una interrupción. La evidencia esperada es baseline, perturbación, métricas, criterio de aborto y conclusión reproducible.

\subsection{Método de análisis}
Ante un caso nuevo, siga esta secuencia:
\begin{enumerate}
\item Identifique actores, dato o recurso compartido y frontera de confianza.
\item Escriba una ejecución nominal y una ejecución adversarial.
\item Distinga seguridad (lo malo no ocurre) de progreso (algo bueno finalmente ocurre).
\item Declare qué supuesto permite avanzar: mayoría, deadline, sincronía parcial, operación idempotente, almacenamiento durable u otro.
\item Compare una alternativa más simple y una más fuerte; registre qué métrica o experimento haría cambiar la decisión.
\end{enumerate}

\subsection{Caso transferible}
Considere una plataforma de reservas que recibe operaciones duplicadas, pierde conectividad entre regiones y debe sostener un camino crítico sin ocultar el estado real al usuario. Aplique el mecanismo de esta clase a una única decisión concreta. Una respuesta madura no intenta resolver todo: delimita la operación, formula la garantía y explica la degradación esperada cuando el supuesto central deja de cumplirse.

\subsection{Errores de razonamiento frecuentes}
\begin{itemize}
\item confundir una propiedad con el producto que suele implementarla;
\item describir el camino feliz sin recorrer una falla intermedia;
\item afirmar ``exactly-once'', ``alta disponibilidad'' o ``consistencia'' sin definir alcance observable;
\item medir solo throughput aunque la hipótesis trate sobre semántica o recuperación;
\item presentar la complejidad añadida como beneficio sin compararla con la alternativa más simple.
\end{itemize}

\section{Lectura guiada}
Principles of Chaos Engineering; Jepsen; literatura de fault injection y testing distribuido.

Durante la lectura, registre: modelo de sistema, propiedad perseguida, contraejemplo que motiva el mecanismo, costo principal y una condición bajo la cual no lo elegiría. Estas cinco anotaciones convierten la bibliografía en una herramienta de diseño y no en una lista para memorizar.

\section{Aplicación al Trabajo Final Integrador}
El equipo debe producir un ADR breve con la decisión asociada a esta clase. Debe contener: contexto, dos alternativas, supuesto decisivo, garantía elegida, trade-off, evidencia pendiente y fecha de revisión. La evidencia mínima esperada para esta unidad es: baseline, perturbación, métricas, criterio de aborto y conclusión reproducible.

\section{Autoevaluación avanzada}
\begin{enumerate}
\item ¿Puedo explicar la clase sin nombrar primero una tecnología?
\item ¿Puedo construir una ejecución que viole la garantía si el mecanismo se configura mal?
\item ¿Puedo distinguir seguridad, progreso y experiencia del usuario?
\item ¿Puedo proponer una medición que valide la propiedad, no solo el rendimiento?
\item ¿Puedo defender cuándo una solución más simple sería suficiente?
\end{enumerate}

\end{document}
